\section{Overlays, logos y texto dinámico}
\label{sec:overlays}

Voctomix 2.0 incorpora un sistema completo de superposición de elementos gráficos
sobre la señal de vídeo: logos corporativos, overlays PNG con canal alfa y texto
dinámico renderizado en tiempo real. Todos estos elementos se gestionan desde la
interfaz gráfica sin interrumpir la retransmisión.

\subsection{Logos}

El sistema soporta logos simultáneos, superpuestos en posiciones fijas del encuadre:

\begin{itemize}
  \item \textbf{logo1}: posicionado en la esquina inferior derecha del fotograma.
  \item \textbf{logo2}: posicionado en la esquina superior izquierda del fotograma.
\end{itemize}

Los ficheros PNG de los logos se definen en las secciones \texttt{[logo1]} y
\texttt{[logo2]} del fichero de configuración, permitiendo adaptar la identidad visual
del evento sin modificar el código. La visibilidad de cada logo se puede activar o
desactivar individualmente desde la interfaz gráfica.

\subsection{Overlays PNG y texto dinámico}

Los overlays son imágenes PNG que se superponen sobre la señal de vídeo
mediante el elemento GStreamer \texttt{gdkpixbufoverlay}. El módulo
\texttt{overlay.py} gestiona el ciclo de vida completo de un overlay:
carga del fichero PNG, fundido de entrada (\textit{blend-in}), mantenimiento en
pantalla y fundido de salida (\textit{blend-out}). El tiempo de fundido se controla
mediante el parámetro \texttt{blend-time} en la sección \texttt{[overlay]} del
fichero de configuración (en milisegundos).

El elemento \texttt{gdkpixbufoverlay} está diseñado para leer archivos de imagen
estáticos desde una ruta del sistema de ficheros. Para cambiar el overlay en caliente,
la clase \texttt{Overlay} actualiza la propiedad \texttt{location} del elemento mediante
el comando \texttt{set\_overlay} del protocolo de control, lo que fuerza la recarga del
nuevo fichero PNG sin interrumpir el pipeline.

\subsection{Rotulación dinámica}
\label{sec:doble_rotulacion}

El sistema de rotulación permite al operador escribir el nombre del ponente o el título
de la sesión en el panel de la interfaz gráfica. Al confirmar el texto, la interfaz lo
envía al núcleo a través del protocolo de control, donde voctocore actualiza
directamente el elemento \texttt{textoverlay} del pipeline GStreamer. El rótulo aparece
en antena en el siguiente fotograma, sin escritura a disco ni procesamiento externo.

El sistema dispone de dos canales de rotulación independientes, lo que permite mostrar
simultáneamente dos textos en posiciones distintas del encuadre sin solapamientos. Cada
canal tiene su propio campo de texto en la interfaz, con gestión de foco de teclado
independiente para evitar colisiones con los atajos del mezclador. Su implementación
requirió extender el pipeline con un segundo elemento \texttt{textoverlay}, ampliar el
protocolo de control con un segundo comando de texto y añadir el campo correspondiente
en la interfaz gráfica.

\subsection{Botón UPDATE: recarga dinámica de recursos}

El botón UPDATE del panel de overlays permite recargar el directorio de
recursos gráficos durante la emisión. Esta funcionalidad es crítica en entornos de
directo: si el equipo de grafismo corrige un error ortográfico en una plantilla PNG o
añade un nuevo rótulo de último minuto, el operador puede incorporarlo al sistema sin
interrumpir la señal de vídeo ni reiniciar ningún proceso.

\subsection{Auto-off de overlays}

La funcionalidad \textit{auto-off} resuelve un problema frecuente en producción en
directo: el operador realiza un corte de cámara olvidando retirar el overlay activo,
que queda superpuesto sobre la nueva fuente de forma inapropiada (por ejemplo, el
nombre de un ponente apareciendo sobre la imagen de otro).

Con esta opción activa, cada vez que el realizador realiza un corte de cámara el
sistema retira automáticamente el overlay y el texto dinámico con un fundido de salida,
sin que el operador tenga que hacerlo manualmente. La funcionalidad puede activarse de
forma permanente en la configuración o exponerse como un botón de alternancia en la
interfaz para que el operador decida en tiempo real si la quiere habilitada.

\begin{figure}[H]
  \centering
  \fbox{\begin{minipage}{0.85\textwidth}\centering\vspace{1.2cm}
    \textit{[IMAGEN PENDIENTE: captura del panel de overlays en voctogui,}\\
    \textit{mostrando los botones LOGO1, LOGO2, AUTO-OFF, UPDATE,}\\
    \textit{los campos de texto dinámico (INSERT 1 e INSERT 2) y el botón de activación.}\\[4pt]
    \textit{Guardar como \texttt{figuras/panel\_overlays.png} y enviar.}
    \vspace{1.2cm}
  \end{minipage}}
  \caption{Panel de overlays y rotulación dinámica en la interfaz de Voctomix~2.0.}
  \label{fig:panel_overlays}
\end{figure}
